\section{Ecological Model and Competing Memory Mechanisms}
\label{sec:model}

\subsection{Strong-Allee predator--prey backbone}
All objects in this subsection are inherited from~\cite{P01}, where they were derived and certified with validated numerics: the model, the locked parameters, the equilibrium $(x^*,y^*(A))$, the Jacobian $J(A)$, and the Matignon boundary $\alpha^*(A)$. They are restated because the experiments of Sections~\ref{sec:numerics}--\ref{sec:benchmark} are defined in terms of them; none is claimed as new.

The nonlinear model is
\begin{align}
 f_1(x,y) &= r x\left(1-\frac{x}{K}\right)\left(\frac{x}{A}-1\right)-\frac{axy}{1+hx},\\
 f_2(x,y) &= e\frac{axy}{1+hx}-my,
\end{align}
with locked parameters
\[
r=\frac32,\qquad K=1,\qquad a=1,\qquad h=\frac12,\qquad e=\frac45,\qquad m=\frac25.
\]
The coexistence equilibrium and Jacobian are
\[
x^*=\frac23,\qquad y^*(A)=\frac{2(2-3A)}{9A},
\]
\[
J(A)=\begin{pmatrix}
\dfrac{7A-2}{8A} & -\dfrac12\\[1ex]
\dfrac{2-3A}{10A} & 0
\end{pmatrix}.
\]
For the commensurate Caputo model, local stability follows Matignon's sector condition $|\arg\mu|>\alpha\pi/2$ on the eigenvalues $\mu$ of $J(A)$~\cite{F05}. In the complex-pair regime $2/7<A<10/19$ the critical order is
\[
\alpha^*(A)=\frac{2}{\pi}\arctan\!\left(\frac{\sqrt{4D(A)-T(A)^2}}{T(A)}\right),
\]
where $T(A)=\operatorname{tr}J$ and $D(A)=\det J$. At $A=0.30$, $\alpha^*\approx0.969012276$: the Caputo system with $\alpha=0.9<\alpha^*$ is stable while its integer-order limit ($\alpha=1>\alpha^*$) is unstable~\cite{P01}. Most discrimination experiments below use $A\le2/7$, where every candidate mechanism is stable at the operating point and stability cannot serve as a discrimination shortcut.

\paragraph{Inherited backbone versus new experimental layer.}
Notation is kept identical to~\cite{P01} for an identification reason: if the ecological equations changed between papers, a difference in simulated observations could reflect the altered biology instead of the memory law. The ecology is therefore frozen and only the response law varies --- ${}^CD_t^\alpha$, the delay $\xi(t-\tau)$, or a latent state $q\in\mathbb{R}^m$ --- with all candidates facing the same budget $\|u\|_\infty\le0.10$, horizon $T=12$, and threshold $A$. Table~\ref{tab:locked-backbone} lists the frozen quantities.

\begin{table}[htbp]
\centering
\caption{Locked ecological backbone used by both manuscripts. These values are inherited inputs, not new contributions of the present paper.}
\label{tab:locked-backbone}
\small
\begin{tabular}{lll}
\hline
\textbf{Quantity} & \textbf{Value} & \textbf{Role in present experiments}\\
\hline
$r$ & $1.5$ & intrinsic prey growth\\
$K$ & $1$ & prey carrying capacity\\
$a$ & $1$ & attack-rate scale\\
$h$ & $0.5$ & Holling-II handling parameter\\
$e$ & $0.8$ & conversion efficiency\\
$m$ & $0.4$ & predator mortality\\
$x^*$ & $2/3$ & common coexistence prey level\\
$y^*(A)$ & $2(2-3A)/(9A)$ & coexistence predator level\\
$J(A)$ & Eq. above & ODE/Caputo local ecological generator\\
\hline
\end{tabular}
\end{table}

The threshold $A$ has two roles. First, $x(t)<A$ defines the failure event: $P(x)<0$ there, and the trajectory can leave the coexistence basin permanently (Figure~\ref{fig:backbone}, left). Second, $x^*-A$ is the natural safety margin; at $A=0.25$ it equals $2/3-0.25\approx0.417$. An input can thus score highly on the information criterion of Section~\ref{sec:design} and still be ecologically inadmissible --- Tables~\ref{tab:benchmark-safety} and~\ref{tab:fd-design-ranking} quantify exactly this.

\begin{figure}[htbp]\centering
\includegraphics[width=0.98\textwidth]{fig11_system_schematic}
\caption{Setting: backbone from~\cite{P01}; perturbation $u(t)$ on the prey channel with $\|u\|_\infty\le0.10$; observations $y_k=Cz(t_k)+\varepsilon_k$ on prey, predator, or both. ODE and Caputo share $J(A)$ exactly; DDE and latent are calibrated rival response laws. Dashed rectangle: numerical benchmark diagnostic region only (not an invariant-set certificate).}
\label{fig:schematic}
\end{figure}

\begin{figure}[htbp]\centering
\includegraphics[width=0.94\textwidth]{fig12_backbone_regime}
\caption{Safety geometry at $A=0.25$ (backbone from~\cite{P01}). Left: recovery for $x_0>A$ versus collapse for $x_0<A$. Right: $P(x)=rx(1-x/K)(x/A-1)$, with sign change at $x=A$.}
\label{fig:backbone}
\end{figure}

\subsection{Four rival mechanisms}
With the reference time $\tau_0$ enforcing dimensional consistency~\cite{F02,F03,F04}, the four candidate mechanisms around the operating point are
\begin{align*}
\text{ODE:}&\quad \dot\xi=J\xi+Bu,\\
\text{Caputo:}&\quad \tau_0^{\alpha-1}{}^CD_t^\alpha\xi=J\xi+Bu,\\
\text{DDE:}&\quad \dot\xi=A_0\xi(t)+A_1\xi(t-\tau)+Bu,\\
\text{latent:}&\quad \dot q=\bar A q+\bar Bu,\qquad y=\bar Cq.
\end{align*}
The Caputo mechanism is the only one carrying a fractional order. The latent class, which would otherwise be flexible enough to imitate anything, is held to an explicit state-dimension budget $m$ --- a restriction whose necessity is precisely the content of Result~II below. In the benchmark the latent \emph{hypothesis} is fixed at $m=3$ while latent \emph{data} are generated at $m=1$ and $m=3$; class-wise statistics for that class are therefore pooled over the two generators, and are labelled \emph{finite-latent} rather than \emph{latent-3} wherever this matters.

\paragraph{Combined fractional-delay stress-test model.}
A system may carry hereditary memory and a trophic lag simultaneously, so we also simulate the combined model
\begin{align}
\tau_0^{\alpha-1}{}^CD_t^\alpha x(t)
&=P(x(t);A)-\frac{a x(t)y(t)}{1+h x(t)}+u(t),\\
\tau_0^{\alpha-1}{}^CD_t^\alpha y(t)
&=e\frac{a x(t-\tau)y(t)}{1+h x(t-\tau)}-m y(t),
\label{eq:fdmodel}
\end{align}
with constant equilibrium history on $[-\tau,0]$. Its boundaries are exact: at $\tau=0$, \eqref{eq:fdmodel} is the pure Caputo model; at $\alpha=1$ it is the retarded DDE benchmark. Both limits are verified numerically to $6.75\times10^{-6}$ and better in Table~\ref{tab:solver-validation}. The model is a numerical stress test only; no theorem depends on it.
